
\chapter{مرور نظام‌مند مجموعه‌داده‌های سی‌تی شکمی}
\label{chap:survey}

فصل~\ref{chap:literature} نشان داد که آموزش مستقیم مدل بخش‌بندی روی داده‌ی بدون‌تزریق می‌تواند دقت این پویش را بالا ببرد، اما این راه به داده‌ی بدون‌تزریقِ برچسب‌خورده‌ی بومی هر مرکز نیاز دارد (بخش~\ref{sec:seg-gap}). این فصل نشان می‌دهد در مجموعه‌داده‌های عمومی سی‌تی شکمی چه مقدار داده‌ی برچسب‌خورده‌ی بدون‌تزریق در دسترس است و این داده‌ها چه محدودیت‌هایی برای اعتبار نتیجه‌ی هر مدل ایجاد می‌کنند.

این موضوع پیش از فصل روش لازم است، چون تعیین می‌کند روش پیشنهادی از چه داده‌ای ساخته شود و با چه داده‌ای ارزیابی شود. مطالب این فصل خلاصه‌ی یک مرور نظام‌مند منتشرشده است~\cite{danaei2026abdomenctreview}؛ تنها یافته‌هایی آمده‌اند که بر طراحی بقیه‌ی کار اثر می‌گذارند و جزئیات کامل روش و جدول استخراج در همان مقاله است. ترتیب فصل چنین است: بخش~\ref{sec:s-protocol} دستورالعمل مرور را تعریف می‌کند، بخش~\ref{sec:s-results} یافته‌ها را به‌همراه آن‌چه از آن‌ها برمی‌آید و دامنه‌ی اعتبارشان می‌آورد، و بخش~\ref{sec:s-implications} پیامد آن‌ها را برای روش پیشنهادی برمی‌شمارد.

\section{دستورالعمل مرور}
\label{sec:s-protocol}

\subsection{انگیزه و محورهای مرور}
\label{sec:datasets}

ادبیات بخش‌بندی سی‌تی شکمی حول شمار محدودی مجموعه‌داده‌ی پایه شکل گرفته است: \lr{LiTS}~\cite{Lits2017}، \lr{KiTS}~\cite{kits2019}، \lr{Pancreas-CT} و \lr{BTCV}. منابع نسل بعد، از جمله \lr{AbdomenCT-1K}~\cite{abdomenct2021}، \lr{FLARE}~\cite{flare23} و \lr{AbdomenAtlas~1.1}~\cite{abdomenatlas2024}، عمدتاً از گردآوری و بازبرچسب‌گذاری همین مجموعه‌ها ساخته شده‌اند. این ساختار درختی دو پیامد شناخته‌شده دارد: نشت داده\footnote{\lr{Data leakage}}، وقتی بیماری هم در آموزش یک پژوهش و هم در ارزیابی پژوهش دیگر حاضر است؛ و تمرکز محک\footnote{\lr{Benchmark concentration}}، یعنی متمرکزشدن استفاده‌ها بر شمار اندکی منبع پرارجاع~\cite{koch2021reduced,varoquaux2022machine}. بنابراین شمار ظاهری مجموعه‌داده‌ها را نمی‌توان بدون ردیابی منشأ، معادل داده‌ی مستقل گرفت.

سوگیری مجموعه‌داده در هوش مصنوعی پزشکی پیشینه‌ی مستقلی دارد~\cite{Obermeyer2019,Larrazabal2020,Seyyed-Kalantari2021,jones2024causal}، اما این مطالعات آن را عمدتاً در \emph{خروجی مدل} می‌سنجند، پس از آن‌که آموزش هر تحریف داده را جذب کرده است~\cite{Kocak2024,Xu2024fairness}. ابزارهای موجود ارزیابی کیفیت نیز برای این منظور طراحی نشده‌اند: \lr{PROBAST+AI}~\cite{Moonse082505} مخاطره\footnote{\lr{Risk}} سوگیری \emph{مدل‌های پیش‌بین} را می‌سنجد و \lr{TRIPOD+AI}~\cite{Collins2024} راهنمای \emph{گزارش‌دهی} است. تا آن‌جا که در جست‌وجوهای این مرور یافت شد، هیچ مرور نظام‌مندی ترکیب جمعیتی، شیوه‌ی برچسب‌گذاری، الگوی بازاستفاده و سهم پویش‌های تزریق را در مجموعه‌داده‌های عمومی سی‌تی شکمی هم‌زمان نسنجیده بود.

بر این پایه، مرور چهار محور را بررسی می‌کند (جدول~\ref{tab:s-questions}). سه محور نخست وضعیت عمومی مجموعه‌داده‌ها را نشان می‌دهند؛ محور چهارم مستقیماً به مسئله‌ی این پایان‌نامه مربوط است، چون ستون پویش جدول استخراج امکان سنجش آن را فراهم می‌کند.

\begin{table}[!htb]
\centering
\footnotesize
\caption{محورهای مرور و آن‌چه در هر محور سنجیده می‌شود.}
\label{tab:s-questions}
\begin{tabular}{|R{0.8cm}|R{4.2cm}|R{6.2cm}|R{1.6cm}|}
\hline
\textbf{شماره} & \textbf{محور} & \textbf{آن‌چه سنجیده می‌شود} & \textbf{بخش} \\
\hline
\عدد{۱} & بازاستفاده و افزونگی & شمار نوبت‌های اصیل در برابر بازاستفاده‌شده، و شبکه‌ی منشأ & \ref{sec:s-reuse} \\
\hline
\عدد{۲} & پوشش جغرافیایی و اندامی & کشورهای مشارکت‌کننده، شمار مجموعه‌های دارای نقاب هر اندام، شیوه‌ی برچسب‌گذاری & \ref{sec:s-geo}، \ref{sec:s-organs} \\
\hline
\عدد{۳} & سوگیری مجموعه‌داده‌ها & امتیاز سه‌سطحی هر مقوله و رده‌ی کلی هر مجموعه & \ref{sec:s-bias} \\
\hline
\عدد{۴} & سهم پویش بدون‌تزریق & سهم نوبت‌های بدون‌تزریق، در سه برآورد & \ref{subsec:s-phase-share} \\
\hline
\end{tabular}
\end{table}

\subsection{جست‌وجو و معیار ورود}

جست‌وجو در \lr{Google Scholar}، \lr{PubMed}، \lr{Scopus} و مخازن نهادی مانند \lr{TCIA} انجام شد. فهرست مراجع و توصیف منبع داده‌ی هر مجموعه نیز بررسی شد تا مجموعه‌های سازنده‌ی بازاستفاده‌شده ردیابی شوند. پنج معیار ورود اعمال شد: حجم سه‌بعدی، دست‌کم یک ناحیه‌ی برچسب‌خورده به‌صورت نقاب، پشتیبانی برچسب‌ها از یک وظیفه‌ی بالینی، مربوط‌بودن برچسب‌ها به ساختارهای شکمی، و استفاده در دست‌کم یک مقاله‌ی داوری‌شده. از \عدد{۴۸} نامزد، \عدد{۴۵} مجموعه‌داده باقی ماند. معیار آخر بدین معناست که مرور، \emph{چشم‌انداز منتشرشده} را نمایندگی می‌کند و نه کل داده‌ی موجود.

\subsection{استخراج داده و ارزیابی سوگیری}

دو ارزیاب مستقل داده را در شش محور استخراج کردند: ترکیب، بافت بالینی، خاستگاه، ویژگی‌های تصویربرداری از جمله پویش تزریق، جزئیات برچسب‌گذاری، و اطلاعات جمعیت‌شناختی. در این فصل، «نوبت» به یک نوبت تصویربرداری سی‌تی (اغلب معادل یک بیمار) به همان معنایی گفته می‌شود که خود مجموعه‌داده گزارش کرده است، و «مورد اصیل» نوبتی است که برای نخستین بار در همان مجموعه منتشر شده و از مجموعه‌ی دیگری وام گرفته نشده است.

چون چک‌لیست پذیرفته‌شده‌ای برای سوگیری مجموعه‌داده‌ها وجود نداشت، یک چارچوب هفت‌مقوله‌ای با تطبیق سه ابزار موجود تدوین شد و روی \عدد{۲۲} مجموعه‌داده‌ای اجرا شد که صد نمونه‌ی عمومی یا بیشتر دارند. هر مقوله در سه سطح مخاطره امتیاز گرفت و امتیازها به یکی از چهار رده‌ی کلی «بحرانی»، «بالا»، «متوسط» یا «پایین» تجمیع شد. جزئیات مقوله‌ها، تعریف رده‌ها، پایایی بین‌ارزیاب و فرمول تجمیع در مقاله‌ی منتشرشده آمده است.

\subsection{تعریف سهم پویش بدون‌تزریق}
\label{sec:s-phase-def}

سنجش سهم پویش بدون‌تزریق به تعریف روشن از ستون پویش نیاز دارد. معیار ورود به مرور وجود نقاب بخش‌بندی بود؛ بنابراین این ستون پویش‌هایی را ثبت می‌کند که نقاب‌های هر مجموعه از آن‌ها می‌آیند و نه لزوماً همه‌ی پویش‌هایی که تصویربرداری شده‌اند. برای نمونه \lr{WAW-TACE}~\cite{sobotka2024wawtace} نوبت‌های چندپویشی و از جمله پویش بدون‌تزریق دارد، اما نقاب‌هایش تنها روی پویش‌های دارای ماده‌ی حاجب ارائه شده‌اند. چون موضوع این محور داده‌ی برچسب‌خورده است، این تعریف با هدف مرور سازگار است.

در سطح مجموعه‌داده شمارش دقیق ممکن است. در سطح شمار نوبت، اما، توصیف‌نامه‌ها تنها حضور هر پویش را اعلام می‌کنند و شمار نوبت‌های هر پویش را نه. از این رو سهم پویش بدون‌تزریق به‌صورت بازه‌ای با سه برآورد تعریف شد: کران پایین قطعی، که تنها نوبت‌های مجموعه‌های تضمین‌شده بدون‌تزریق را می‌شمارد؛ کران بالا، که همه‌ی نوبت‌های مجموعه‌هایی را می‌شمارد که این پویش را پوشش می‌دهند و آشکارا بیش‌برآورد است؛ و برآورد میانی، که نوبت‌های هر مجموعه را یکنواخت میان پویش‌هایش پخش می‌کند.

\section{یافته‌ها و نتیجه‌گیری}
\label{sec:s-results}


\subsection{بازاستفاده و افزونگی}
\label{sec:s-reuse}

از \عدد{۴۵} مجموعه‌داده، مجموعاً \عدد{۴۷٬۰۴۹} نوبت سی‌تی فهرست شده است و تنها \عدد{۲۳٬۷۴۶} مورد آن \emph{اصیل} بود؛ یعنی \عدد{۴۹٫۵} درصد از حجم ظاهری بازاستفاده‌ی موارد پیشین است و شش مجموعه‌داده اصلاً مورد اصیلی ندارند. انتشار مجموعه‌داده‌ها از سال \عدد{۲۰۲۲} جهش کرد و در \عدد{۲۰۲۳} با \عدد{۱۱} مجموعه به اوج رسید؛ در همان سال کمتر از نیمی از نوبت‌های منتشرشده اصیل بودند.

شکل~\ref{fig:review-reuse} منشأ این بازاستفاده را نشان می‌دهد. چند مجموعه‌ی پایه گره‌های مرکزی شبکه‌اند و مواردشان به مجموعه‌های ترکیبی متأخرتر سرریز می‌شود؛ بزرگ‌ترین آن‌ها، \lr{AbdomenAtlas~1.1}، حدود \عدد{۶۰} درصد مواردش را از مجموعه‌های دیگر گرفته است.

\begin{figure}[H]
\centering
\includegraphics[width=0.9\textwidth]{review/dataset_reuse_network.png}
\caption{شبکه‌ی بازاستفاده‌ی موارد میان مجموعه‌داده‌های عمومی سی‌تی شکمی. هر گره یک مجموعه‌داده است و اندازه‌ی آن با شمار سری‌های سی‌تی متناسب است؛ آبی یعنی تماماً اصیل، نارنجی ترکیبی از اصیل و بازاستفاده، و قرمز تماماً بازاستفاده. هر پیکان از مجموعه‌ی منبع به مجموعه‌ی مشتق می‌رود. برگرفته از~\cite{danaei2026abdomenctreview}.}
\label{fig:review-reuse}
\end{figure}

شمار مجموعه‌داده‌ها داده‌ی مستقل را بیش‌برآورد می‌کند: با بازاستفاده‌ی \عدد{۴۹٫۵} درصدی، حجم ظاهری این زیست‌بوم تقریباً دو برابر حجم داده‌ی اصیل آن است. ارزیابی روی مجموعه‌های ترکیبی از این رو می‌تواند بیمارانی را شامل شود که در آموزش پژوهش دیگر نیز حاضر بوده‌اند، مگر آن‌که منشأ هر مورد ردیابی شود.

\subsection{پوشش جغرافیایی}
\label{sec:s-geo}

دست‌کم \عدد{۲۴} کشور در مجموعه‌داده‌ها مشارکت دارند، اما \عدد{۷۵٫۳} درصد مجموعه‌ها از ایالات متحده، کانادا و کشورهای اروپایی می‌آیند. شکل~\ref{fig:review-geo} نشان می‌دهد که آفریقای زیرصحرا، جنوب آسیا و بخش عمده‌ی خاورمیانه هیچ نمایندگی‌ای ندارند.

\begin{figure}[H]
\centering
\includegraphics[width=\textwidth]{review/geographic_distribution.png}
\caption{توزیع جغرافیایی کشورهای مشارکت‌کننده در مجموعه‌داده‌های عمومی سی‌تی شکمی. رنگ هر کشور شمار مجموعه‌داده‌هایی است که در آن‌ها سهم دارد و کشورهای خاکستری سهمی ندارند. برگرفته از~\cite{danaei2026abdomenctreview}.}
\label{fig:review-geo}
\end{figure}

نبودِ نمایندگی جغرافیایی اعتبار نتایج را محدود می‌کند: مدلی که تنها داده‌ی چند منطقه را دیده باشد، برای جمعیت‌های نمایندگی‌نشده شواهدی از اعتبار ندارد.

\subsection{پوشش اندام و شیوه‌ی برچسب‌گذاری}
\label{sec:s-organs}

کبد در \عدد{۲۳} مجموعه‌داده، کلیه و پانکراس هر یک در \عدد{۱۹} و طحال در \عدد{۱۴} مجموعه نقاب دارند. در مقابل، عروق کبدی تنها در \عدد{۵} مجموعه و تنه‌ی سلیاک و عروق مزانتریک هر یک در \عدد{۲} مجموعه برچسب خورده‌اند.

از \عدد{۴۵} مجموعه‌داده، \عدد{۲۷} مورد (\عدد{۶۰} درصد) تماماً به‌دست متخصص و \عدد{۱۷} مورد (\عدد{۳۷٫۸} درصد) با زنجیره‌ای نیمه‌خودکار\footnote{\lr{Semi-automatic}} برچسب خورده‌اند که خروجی یک مدل را متخصص اصلاح می‌کند. زنجیره‌ی نیمه‌خودکار در شمار مجموعه‌ها کمتر از نیمی است، اما بزرگ‌ترین مجموعه‌ها (\lr{AbdomenAtlas~1.1}، \lr{FLARE}، \lr{TotalSegmentator} و \lr{AMOS22}) همگی در این گروه‌اند و از این رو در حجم داده غالب است.

پوشش اندام‌ها یکنواخت نیست: اندام‌های بزرگ در بسیاری از مجموعه‌ها برچسب دارند و ساختارهای عروقی از کم‌برچسب‌ترین ساختارها هستند. چون بزرگ‌ترین مجموعه‌ها با زنجیره‌ی نیمه‌خودکار برچسب خورده‌اند، کیفیت نقاب‌های مرجع در بخش غالب حجم داده به کیفیت مدلی وابسته است که در آن زنجیره خروجی اولیه را می‌سازد.

\subsection{مخاطره‌ی سوگیری}
\label{sec:s-bias}

شکل~\ref{fig:review-bias} توزیع مخاطره را در هفت مقوله نشان می‌دهد. سوگیری نژادی\footnote{\lr{Racial bias}} (\عدد{۱۶} از \عدد{۲۲}) و سوگیری انتخاب\footnote{\lr{Selection bias}} (\عدد{۱۵} از \عدد{۲۲}) پرتکرارترین مقوله‌های پرمخاطره‌اند: اغلب مجموعه‌ها تک‌مرکزی‌اند یا ترکیب نژادی را گزارش نمی‌کنند، و گروه شاهد آن‌ها اغلب فقط موارد سالم است. سوگیری برچسب‌گذاری تنها در \عدد{۵} مجموعه پرمخاطره است. سوگیری جمعیت‌شناختی و زمانی هیچ مورد پرمخاطره‌ای ندارند، اما بیشترین سهم «در دسترس نیست» در همین دو مقوله است (\عدد{۶} و \عدد{۸} مجموعه).

\begin{figure}[H]
\centering
\includegraphics[width=\textwidth]{review/bias_risk_distribution.png}
\caption{توزیع سطح مخاطره در هفت مقوله‌ی سوگیری برای \عدد{۲۲} مجموعه‌داده‌ی دارای صد نمونه‌ی عمومی یا بیشتر. قرمز مخاطره‌ی بالا، زرد متوسط، سبز پایین و خاکستری «در دسترس نیست» است. برگرفته از~\cite{danaei2026abdomenctreview}.}
\label{fig:review-bias}
\end{figure}

جدول~\ref{tab:review-overall} رده‌ی کلی هر مجموعه را نشان می‌دهد: \عدد{۱۲} مجموعه از \عدد{۲۲} مجموعه در رده‌ی بحرانی‌اند و تنها \عدد{۳} مورد در رده‌ی پایین. دو مورد از این سه، \lr{AbdomenAtlas~1.1}~\cite{abdomenatlas2024} با \عدد{۱۱۲} مرکز و \lr{RSNA Abdominal Trauma}~\cite{rudie2024rsna} با \عدد{۲۳} مرکز، چندمرکزی و بزرگ‌مقیاس‌اند. در مقابل، \lr{WORD}~\cite{word2022}، تنها مجموعه‌ی اختصاصی بدون‌تزریق، تک‌مرکزی است و در رده‌ی بحرانی قرار دارد.

\begin{table}[!htb]
\centering
\footnotesize
\caption{دسته‌بندی کلی سوگیری \عدد{۲۲} مجموعه‌داده‌ی ارزیابی‌شده.}
\label{tab:review-overall}
\begin{tabular}{|R{2.2cm}|R{1.2cm}|R{9.4cm}|}
\hline
\textbf{رده‌ی سوگیری} & \textbf{شمار} & \textbf{مجموعه‌داده‌ها} \\
\hline
بحرانی & \عدد{۱۲} & {\lr{\mbox{AutoPET}, \mbox{CARE}, \mbox{Colorectal-Liver-Metastases}, \mbox{HCC-TACE-Seg}, \mbox{KiPA22}, \mbox{KiTS}, \mbox{MSD-Colon}, \mbox{MSD-HepaticVessel}, \mbox{MSD-Pancreas}, \mbox{Pancreatic Cancer Survival Prediction}, \mbox{StageII-Colorectal-CT}, \mbox{WORD}}\par} \\
\hline
بالا & \عدد{۲} & {\lr{\mbox{WAW-TACE}, \mbox{WORC-GIST}}\par} \\
\hline
متوسط & \عدد{۵} & {\lr{\mbox{AMOS22}, \mbox{LiTS}, \mbox{PANORAMA}, \mbox{TotalSegmentator}, \mbox{ULS23}}\par} \\
\hline
پایین & \عدد{۳} & {\lr{\mbox{AbdomenAtlas 1.1}, \mbox{RSNA Abdominal Trauma}, \mbox{Pediatric-CT-SEG}}\par} \\
\hline
\end{tabular}
\end{table}

در \عدد{۲۲} مجموعه‌ی ارزیابی‌شده، ضعف غالب در ترکیب داده است و نه در کیفیت حاشیه‌نویسی: سوگیری نژادی و انتخاب پرتکرارترند و سوگیری برچسب‌گذاری کم‌تکرار است. نبودِ مورد پرمخاطره در سوگیری جمعیت‌شناختی و زمانی را نباید کم‌مخاطره‌بودن خواند، چون بیشترین سهم «در دسترس نیست» در همین دو مقوله است، یعنی سن، جنس یا بازه‌ی گردآوری اصلاً گزارش نشده است. سه مجموعه‌ی رده‌ی پایین با نقش تنوع مراکز سازگارند، اما با سه مورد نمی‌توان اثر تنوع مراکز را از اثر اندازه‌ی مجموعه جدا کرد.

\subsection{سهم پویش بدون‌تزریق}
\label{subsec:s-phase-share}

در سطح مجموعه‌داده، از \عدد{۴۵} مجموعه، \عدد{۳۲} مورد (\عدد{۷۱٫۱} درصد) نقاب‌های خود را تنها از پویش‌های دارای ماده‌ی حاجب می‌گیرند. تنها \عدد{۱۲} مورد (\عدد{۲۶٫۷} درصد) پویش بدون‌تزریق را در میان پویش‌های حاشیه‌نویسی‌شده دارند، و از این \عدد{۱۲} مورد تنها \lr{WORD} اختصاصاً بدون‌تزریق است (\عدد{۱۷۰} نوبت، که \عدد{۱۵۰} مورد آن اصیل است)؛ بقیه مجموعه‌های ترکیبی‌اند. یک مجموعه هیچ اطلاعات پویشی گزارش نمی‌کند.

در سطح شمار نوبت، جدول~\ref{tab:review-phase} سه برآورد تعریف‌شده در بخش~\ref{sec:s-phase-def} را کنار هم می‌گذارد. کران پایین قطعی \عدد{۰٫۳} درصد و کران بالا \عدد{۷۰٫۵} درصد است و دو برآورد میانی حدود \عدد{۲۲} درصد به دست می‌دهند. برآورد میانی معادل نسبت نوبت‌های با ماده‌ی حاجب به بدون‌تزریق حدود \عدد{۳٫۴} به \عدد{۱} است.

\begin{table}[!htb]
\centering
\footnotesize
\caption{برآورد سهم نوبت‌های بدون‌تزریق در داده‌ی برچسب‌خورده‌ی عمومی سی‌تی شکمی، محاسبه‌شده از جدول استخراج مرور~\cite{danaei2026abdomenctreview}.}
\label{tab:review-phase}
\begin{tabular}{|R{5.4cm}|R{2.4cm}|R{2.2cm}|R{1.8cm}|}
\hline
\textbf{برآورد} & \textbf{نوبت بدون‌تزریق} & \textbf{از مجموع} & \textbf{سهم (درصد)} \\
\hline
کران پایین قطعی (تنها \lr{WORD}) & \عدد{۱۵۰} & \عدد{۴۷٬۰۴۹} & \عدد{۰٫۳} \\
\hline
برآورد میانی، پخش یکنواخت & \عدد{۹٬۸۵۰} & \عدد{۴۳٬۷۳۶} & \عدد{۲۲٫۵} \\
\hline
برآورد میانی، تنها بخش عمومی & \عدد{۶٬۴۹۳} & \عدد{۲۹٬۸۷۰} & \عدد{۲۱٫۷} \\
\hline
کران بالا & \عدد{۳۳٬۱۸۳} & \عدد{۴۷٬۰۴۹} & \عدد{۷۰٫۵} \\
\hline
\end{tabular}
\end{table}

در دو برآورد میانیِ مبتنی بر پخش یکنواخت، سهم داده‌ی برچسب‌خورده‌ی بدون‌تزریق حدود \عدد{۲۲} درصد است. این عدد برآوردی تحت فرض توزیع فازهاست و نباید با کران بالای \عدد{۷۰٫۵} درصدی جدول یکسان گرفته شود. فاصله‌ی میان کران‌ها نیز نشان می‌دهد که گزارش حضور یک فاز، جای شمارش دقیق موارد آن را نمی‌گیرد. نتیجه‌ی کاربردی برای این پایان‌نامه، محدودیت و نامعلومی پوشش داده‌ی بدون‌تزریق است، نه تعیین یک سقف قطعی برای آن.

\subsection{دامنه‌ی اعتبار}
\label{sec:s-interpretation}

هر یافته‌ی بالا در محدوده‌ی سه قید تفسیر می‌شود:

\begin{itemize}
\item مرور تنها مجموعه‌داده‌هایی را می‌شمارد که دست‌کم در یک مقاله‌ی داوری‌شده به کار رفته‌اند؛ بنابراین یافته‌ها به \emph{چشم‌انداز منتشرشده} مربوط‌اند و داده‌ی بومی بیمارستان‌ها را نمی‌پوشانند.
\item ارزیابی سوگیری تنها روی \عدد{۲۲} مجموعه‌ی دارای صد نمونه‌ی عمومی یا بیشتر انجام شد؛ پس نتیجه‌ی آن درباره‌ی مجموعه‌های بزرگ‌تر است و به مجموعه‌های کوچک‌تر تعمیم داده نمی‌شود.
\item سهم بدون‌تزریق بر ستون پویشِ نقاب‌ها و بر شمارشی متکی است که واحدش همان «نوبت» گزارش‌شده‌ی خود مجموعه‌هاست؛ از این رو برآوردهای بازه‌ای تنها برای مقایسه‌ی مرتبه‌ی بزرگی به کار می‌روند.
\end{itemize}

\section{پیامد برای روش پیشنهادی}
\label{sec:s-implications}

این یافته‌ها چهار پیامد مشخص دارند که هر یک بر طراحی بقیه‌ی کار اثر می‌گذارد.

یکم، آموزش مستقیم روی پویش بدون‌تزریق یک راه معتبر است، اما دسترسی به داده‌ی مناسب آن یکنواخت نیست: حدود \عدد{۷۱} درصد مجموعه‌داده‌ها نقاب بدون‌تزریق ندارند و سهم این پویش در برآورد میانی حدود یک‌پنجم است. پارک و همکاران~\cite{park2025nect} امکان عملکرد مطلوب با آموزش روی این پویش را نشان دادند. یافته‌های مرور، بررسی راهی مکمل را توجیه می‌کنند که از داده‌ی چندپویشی موجود بهره بگیرد؛ کفایت آموزش مستقیم به شمار نمونه‌ها، اندام هدف و کیفیت برچسب‌ها نیز وابسته است.

دوم، این پایان‌نامه به‌جای بازآموزی بخش‌بند، ورودی آن را با هم‌نهادی تصویر دارای ماده‌ی حاجب از تصویر بدون‌تزریق تغییر می‌دهد. هسته‌ی این مسیر دو گام دارد: هم‌ترازسازی پویش‌های موجود یک بیمار برای ساخت جفت‌های آموزشی، و یادگیری تبدیل تصویر بدون‌تزریق به تصویر دارای ماده‌ی حاجب. بنابراین خروجی این مسیر، تصویر هم‌نهاده برای یک بخش‌بند ثابت است، نه مجموعه‌ای تازه با حاشیه‌نویسی انسانی. تشخیص فاز هر سری نیز ابزار جانبی آماده‌سازی و کنترل داده برای این دو گام است.

سوم، ارزیابی روی یک مجموعه‌داده‌ی اصیل و منفرد انجام می‌شود. بازاستفاده‌ی \عدد{۴۹٫۵} درصدی ارزیابی روی مجموعه‌های ترکیبی را با خطر نشت داده همراه می‌کند؛ به همین دلیل آزمایش‌های این پایان‌نامه روی \lr{VinDr-Multiphase}~\cite{zhou2021phaserecog}، مجموعه‌ای اصیل و عمومی، انجام می‌شود و تقسیم داده در سطح بیمار و روی یک پرونده\footnote{\lr{File}} منجمد است.

چهارم، ارقام آزمایش‌ها باید در دامنه‌ی اعتبار مرور تفسیر شوند. سوگیری جغرافیایی و جمعیت‌شناختی مستندشده، و وابستگی نقاب‌های مرجع به زنجیره‌ی نیمه‌خودکار برچسب‌گذاری، دامنه‌ی اعتبار ارقام فصل~\ref{chap:experiments} را محدود می‌کنند و هنگام تفسیر آن‌ها در نظر گرفته می‌شوند.

\section{جمع‌بندی}
\label{sec:s-summary}

این فصل نشان داد که در مجموعه‌داده‌های عمومی سی‌تی شکمی حدود نیمی از نوبت‌ها بازاستفاده‌اند، ترکیب داده از نظر جغرافیا و جمعیت نامتوازن است، و سهم داده‌ی برچسب‌خورده‌ی بدون‌تزریق در برآورد میانیِ فرض‌محور حدود یک‌پنجم است. این یافته‌ها انگیزه‌ی استفاده از داده‌ی چندپویشی موجود را فراهم می‌کنند و بر اهمیت تفکیک داده و ارزیابی اندام‌محور تأکید دارند. فصل~\ref{chap:method} مسیر عملی این انتخاب را معرفی می‌کند: ابزار تشخیص فاز، زنجیره‌ی هم‌ترازسازی برای آماده‌سازی جفت‌ها، و شبکه‌ای که تصویر دارای ماده‌ی حاجب را از تصویر بدون‌تزریق هم‌نهاده می‌کند.
